Vision-language models (VLMs) are increasingly proposed for chest X-ray (CXR) interpretation, yet their reliability under distribution shift remains poorly characterised. We study whether CXR VLMs can be selectively trusted when image quality and acquisition conditions change, and whether apparent reliability is grounded in visual evidence. We evaluate MedGemma-4B-IT and Qwen2.5-VL-7B-Instruct on a 1,000-study NIH ChestX-ray14 pilot, a six-axis image-quality shift benchmark, counterfactual image-reliance tests, and a 264-sample CheXpert-plus-RRG external report set. Both models produce structured JSON predictions with zero parse failures, but their reliability profiles diverge sharply. MedGemma is strongly image-dependent (image sensitivity: 1.00 blank, 0.95 blur, 0.88 shuffle) yet only moderately discriminative (macro AUROC 0.62--0.65) and poorly calibrated on held-out data (macro ECE 0.22). Qwen2.5-VL is effectively image-invariant (image sensitivity: 0.003 across blank, blur, and shuffle), predicts \texttt{no\_finding} for 988/1,000 studies, and remains at chance discrimination (macro AUROC $\approx$ 0.50), while still producing valid JSON and apparently favourable selective metrics. Shift effects are heterogeneous and non-monotonic: several blur, noise, JPEG, and resolution conditions do not degrade, and sometimes improve, macro AUROC or ECE. On CheXpert-plus-RRG, MedGemma shows 13.5\% report-level hallucination and 12.2\% omission, while original-vs-blank status flips occur in 45.3\% of pairs. These findings show that accuracy, calibration, or selective risk alone can mischaracterise CXR VLM reliability. Reliable selective trust requires joint evaluation of calibration, uncertainty, shift response, and visual-evidence reliance.
